\section{Conclusion}

We presented the first systematic comparison of token entropy and semantic consistency for hallucination detection. Our key findings:

\begin{enumerate}
    \item \textbf{Consistency dramatically outperforms entropy:} AUROC 0.81 versus 0.65, Cohen's $d$ 1.19 versus 0.47
    \item \textbf{Linear fusion provides no improvement} at proof-of-concept scale ($N$=20)
    \item \textbf{Signals are moderately correlated} ($r$=-0.54), suggesting partial redundancy
\end{enumerate}

These results challenge the assumption that combining uncertainty signals automatically improves detection. For practitioners building hallucination detection systems, consistency should be the foundation. The complementarity hypothesis requires larger-scale validation before practical deployment.

\textbf{Future work} should validate on larger samples ($N \geq 500$), extend to additional benchmarks (Natural Questions, TruthfulQA), explore non-linear combination strategies, and test across model scales (13B, 70B).
